\documentclass[10pt,letterpaper]{article}
\usepackage[utf8]{inputenc}
\usepackage[intlimits]{amsmath}
\usepackage{amsfonts}
\usepackage{amssymb}
\usepackage{ragged2e}
\usepackage[letterpaper, margin=0.5in]{geometry}
\usepackage{graphicx}
\usepackage{cancel}
\usepackage{mathtools}
\usepackage{tabularx}
\usepackage{arydshln}
\usepackage{tensor}
\usepackage{array}
\usepackage{xcolor}
\usepackage[boxed]{algorithm}
\usepackage[noend]{algpseudocode}
\usepackage{listings}
\usepackage{textcomp}
% \usepackage[pdf,tmpdir,singlefile]{graphviz}
\usepackage{mathrsfs}
\usepackage{bbm}
\usepackage{tikz}
\usepackage{tikz-cd}
\usepackage{enumitem}
\usepackage{arydshln}
\usepackage{relsize}
\usepackage{multicol}
\usepackage{scalerel}
\usepackage{upgreek}
\usepackage{wasysym}
\usepackage{ifthen}

\usetikzlibrary{bayesnet}
\setlist{noitemsep}

%%%%%%%%%%%%%%%%%%%%%%%%%%%%%
% Formatting commands
%%%%%%%%%%%%%%%%%%%%%%%%%%%%%
\newcommand{\n}{\hfill\break}
\newcommand{\nn}{\vspace{0.5\baselineskip}\n}
\newcommand{\up}{\vspace{-\baselineskip}}
\newcommand{\hangblock}[2]{\par\noindent\settowidth{\hangindent}{\textbf{#1: }}\textbf{#1: }\nolinebreak#2}
\newcommand{\lemma}[1]{\hangblock{Lemma}{#1}}
\newcommand{\defn}[1]{\hangblock{Defn}{#1}}
\newcommand{\thm}[1]{\hangblock{Thm}{#1}}
\newcommand{\cor}[1]{\hangblock{Cor}{#1}}
\newcommand{\prop}[1]{\hangblock{Prop}{#1}}
\newcommand{\ex}[1]{\hangblock{Ex}{#1}}
\newcommand{\exer}[1]{\hangblock{Exer}{#1}}
\newcommand{\fact}[1]{\hangblock{Fact}{#1}}
\newcommand{\remark}[1]{\hangblock{Remark}{#1}}
\newcommand{\proven}{\;$\square$\n}
\newcommand{\problem}[1]{\par\noindent{\nolinebreak#1}\n}
\newcommand{\problempart}[2]{\par\noindent\indent{}\settowidth{\hangindent}{\textbf{(#1)} \indent{}}\textbf{(#1) }\nolinebreak#2\n}
\newcommand{\ptxt}[1]{\textrm{\textnormal{#1}}}
\newcommand{\inlineeq}[1]{\centerline{$\displaystyle #1$}}
\newcommand{\pageline}{\up\par\noindent\rule{\textwidth}{0.1pt}}

%%%%%%%%%%%%%%%%%%%%%%%%%%%%%
% Math commands
%%%%%%%%%%%%%%%%%%%%%%%%%%%%%
% Set Theory
\newcommand{\card}[1]{\left|#1\right|}
\newcommand{\set}[1]{\left\{#1\right\}}
\newcommand{\setmid}{\;\middle|\;}
\newcommand{\ps}[1]{\mathcal{P}\left(#1\right)}
\newcommand{\pfinite}[1]{\mathcal{P}^{\ptxt{finite}}\left(#1\right)}
\newcommand{\naturals}{\mathbb{N}}
\newcommand{\N}{\naturals}
\newcommand{\integers}{\mathbb{Z}}
\newcommand{\Z}{\integers}
\newcommand{\rationals}{\mathbb{Q}}
\newcommand{\Q}{\rationals}
\newcommand{\reals}{\mathbb{R}}
\newcommand{\R}{\reals}
\newcommand{\complex}{\mathbb{C}}
\newcommand{\C}{\complex}
\newcommand{\halfPlane}{\mathbb{H}}
\let\HSym\H
\let\H\relax
\newcommand{\H}{\halfPlane}
\newcommand{\comp}{^{\complement}}
\DeclareMathOperator{\Hom}{Hom}
\newcommand{\Ind}{\mathbbm{1}}
\newcommand{\cut}{\setminus}
\DeclareMathOperator{\elem}{elem}

% Differentiable Manifolds
\newcommand{\RP}{\mathbb{RP}}
\newcommand{\CP}{\mathbb{RP}}
\newcommand{\osubset}{\overset{\mathclap{\scalebox{0.5}{\ptxt{open}}}}{\subset}}
\newcommand{\osubseteq}{\overset{\mathclap{\scalebox{0.5}{\ptxt{open}}}}{\subseteq}}
\newcommand{\osupset}{\overset{\mathclap{\scalebox{0.5}{\ptxt{open}}}}{\supset}}
\newcommand{\osupseteq}{\overset{\mathclap{\scalebox{0.5}{\ptxt{open}}}}{\supseteq}}
\newcommand{\pdat}[3]{\left.\pd{#1}{#2}\right|_{#3}}
\DeclareMathOperator{\so}{so}
\DeclareMathOperator{\codim}{codim}
\DeclareMathOperator{\Diff}{Diff}
\let\dSym\d
\let\d\relax
\newcommand{\d}{\partial}
\DeclareMathOperator{\gl}{gl}

% Graph Theory
\let\deg\relax
\DeclareMathOperator{\deg}{deg}
\newcommand{\degp}{\ptxt{deg}^{+}}
\newcommand{\degn}{\ptxt{deg}^{-}}
\newcommand{\precdot}{\mathrel{\ooalign{$\prec$\cr\hidewidth\hbox{$\cdot\mkern0.5mu$}\cr}}}
\newcommand{\succdot}{\mathrel{\ooalign{$\cdot\mkern0.5mu$\cr\hidewidth\hbox{$\succ$}\cr\phantom{$\succ$}}}}
\DeclareMathOperator{\cl}{cl}
\DeclareMathOperator{\affdim}{affdim}

% Probability
\newcommand{\parSymbol}{\P}
\newcommand{\Prob}{\mathbb{P}}
\renewcommand{\P}{\Prob}
\newcommand{\Avg}{\mathbb{E}}
\newcommand{\E}{\Avg}
\DeclareMathOperator{\Var}{Var}
\DeclareMathOperator{\cov}{cov}
\DeclareMathOperator{\Unif}{Unif}
\DeclareMathOperator{\Binom}{Binom}
\newcommand{\CI}{\mathrel{\text{\scalebox{1.07}{$\perp\mkern-10mu\perp$}}}}

% Standard Math
\newcommand{\inv}{^{-1}}
\newcommand{\abs}[1]{\left|#1\right|}
\newcommand{\ceil}[1]{\left\lceil{}#1\right\rceil{}}
\newcommand{\floor}[1]{\left\lfloor{}#1\right\rfloor{}}
\newcommand{\conj}[1]{\overline{#1}}
\newcommand{\of}{\circ}
\newcommand{\tri}{\triangle}
\newcommand{\inj}{\hookrightarrow}
\newcommand{\surj}{\twoheadrightarrow}
\newcommand{\ndiv}{\nmid}
\renewcommand{\epsilon}{\varepsilon}
\newcommand{\divides}{\mid}
\newcommand{\ndivides}{\nmid}
\DeclareMathOperator{\lcm}{lcm}
\DeclareMathOperator{\sgn}{sgn}
\newcommand{\map}[4]{\!\!\!\begin{array}[t]{rcl}#1 & \!\!\!\!\to & \!\!\!\!#2\\ {}#3 & \!\!\!\!\mapsto & \!\!\!\!#4\end{array}}
\newcommand{\bigsum}[2]{\smashoperator[lr]{\sum_{\scalebox{#1}{$#2$}}}}
\DeclareMathOperator{\gcf}{gcf}
\newcommand{\restr}[2]{\left.#1\right|_{#2}}

% Linear Algebra
\newcommand{\Id}{\textrm{\textnormal{Id}}}
\newcommand{\im}{\textrm{\textnormal{im}}}
\newcommand{\norm}[1]{\abs{\abs{#1}}}
\newcommand{\tpose}{^{T}\!}
\newcommand{\iprod}[1]{\left<#1\right>}
\newcommand{\giprod}{\iprod{\;\,,\;}}
\DeclareMathOperator{\tr}{tr}
\DeclareMathOperator{\trace}{tr}
\newcommand{\chgBasMat}[3]{\!\!\tensor*[_{#1}]{\left[#2\right]}{_{#3}}}
\newcommand{\vecBas}[2]{\tensor*[]{\left[#1\right]}{_{#2}}}
\DeclareMathOperator{\GL}{GL}
\DeclareMathOperator{\Mat}{Mat}
\DeclareMathOperator{\vspan}{span}
\DeclareMathOperator{\rank}{rank}
\newcommand{\V}[1]{\vec{#1}}
\DeclareMathOperator{\proj}{proj}
\DeclareMathOperator{\compProj}{comp}
\DeclareMathOperator{\row}{row}
\newcommand{\smallPMatrix}[1]{\paren{\begin{smallmatrix}#1\end{smallmatrix}}}
\newcommand{\smallBMatrix}[1]{\brack{\begin{smallmatrix}#1\end{smallmatrix}}}
\newcommand{\pmat}[1]{\begin{pmatrix}#1\end{pmatrix}}
\newcommand{\bmat}[1]{\begin{bmatrix}#1\end{bmatrix}}
\newcommand{\dual}{^{*}}
\newcommand{\pinv}{^{\dagger}}
\newcommand{\horizontalMatrixLine}{\ptxt{\rotatebox[origin=c]{-90}{$|$}}}
\DeclareMathOperator{\range}{range}
\DeclareMathOperator{\Symm}{Symm}
\DeclareMathOperator{\SU}{SU}

% Multilinear Algebra
\let\Lsym\L
\let\L\relax
\DeclareMathOperator{\L}{\mathscr{L}}
\DeclareMathOperator{\A}{\mathcal{A}}
\DeclareMathOperator{\Alt}{Alt}
\DeclareMathOperator{\Sym}{Sym}
\newcommand{\ot}{\otimes}
\newcommand{\ox}{\otimes}
\DeclareMathOperator{\asc}{asc}
\DeclareMathOperator{\asSet}{set}
\DeclareMathOperator{\sort}{sort}
\DeclareMathOperator{\ringA}{\mathring{A}}
\DeclareMathOperator{\Sh}{Sh}

% Topology
\newcommand{\closure}[1]{\overline{#1}}
\newcommand{\uball}{\mathcal{U}}
\DeclareMathOperator{\Int}{Int}
\DeclareMathOperator{\Ext}{Ext}
\DeclareMathOperator{\Bd}{Bd}
\DeclareMathOperator{\rInt}{rInt}
\DeclareMathOperator{\ch}{ch}
\DeclareMathOperator{\ah}{ah}
\newcommand{\LargerTau}{\mathlarger{\mathlarger{\mathlarger{\mathlarger{\tau}}}}}
\newcommand{\Tau}{\mathcal{T}}

% Analysis
\DeclareMathOperator{\Graph}{Graph}
\DeclareMathOperator{\epi}{epi}
\DeclareMathOperator{\hypo}{hypo}
\DeclareMathOperator{\supp}{supp}
\newcommand{\lint}[2]{\underset{#1}{\overset{#2}{{\color{black}\underline{{\color{white}\overline{{\color{black}\int}}\color{black}}}}}}}
\newcommand{\uint}[2]{\underset{#1}{\overset{#2}{{\color{white}\underline{{\color{black}\overline{{\color{black}\int}}\color{black}}}}}}}
\newcommand{\alignint}[2]{\underset{#1}{\overset{#2}{{\color{white}\underline{{\color{white}\overline{{\color{black}\int}}\color{black}}}}}}}
\newcommand{\extint}{\ptxt{ext}\int}
\newcommand{\extalignint}[2]{\ptxt{ext}\alignint{#1}{#2}}
\newcommand{\conv}{\ast}
\newcommand{\pd}[2]{\frac{\partial{}#1}{\partial{}#2}}
\newcommand{\del}{\nabla}
\DeclareMathOperator{\grad}{grad}
\DeclareMathOperator{\curl}{curl}
\let\div\relax
\DeclareMathOperator{\div}{div}
\DeclareMathOperator{\vol}{vol}

% Complex Analysis
\let\Re\relax
\DeclareMathOperator{\Re}{Re}
\let\Im\relax
\DeclareMathOperator{\Im}{Im}
\DeclareMathOperator{\Res}{Res}

% Abstract Algebra
\DeclareMathOperator{\ord}{ord}
\newcommand{\generated}[1]{\left<#1\right>}
\newcommand{\cycle}[1]{\smallPMatrix{#1}}
\newcommand{\id}{\ptxt{id}}
\newcommand{\iso}{\cong}
\DeclareMathOperator{\Aut}{Aut}
\DeclareMathOperator{\SL}{SL}
\DeclareMathOperator{\op}{op}
\newcommand{\isom}[4]{\!\!\!\begin{array}[t]{rcl}#1 & \!\!\!\!\overset{\sim}{\to} & \!\!\!\!#2\\ #3 & \!\!\!\!\mapsto & \!\!\!\!#4\end{array}}
\newcommand{\F}{\mathbb{F}}
\newcommand{\acts}{\;\reflectbox{\rotatebox[origin=c]{-90}{$\circlearrowright$}}\;}
\newcommand{\disjunion}{\mathrel{\text{\scalebox{1.07}{$\perp\mkern-10mu\perp$}}}}
\DeclareMathOperator{\SO}{SO}
\DeclareMathOperator{\stab}{stab}
\DeclareMathOperator{\nullity}{nullity}
\DeclareMathOperator{\Perm}{Perm}
\DeclareMathOperator{\nsubgp}{\vartriangleleft}
\DeclareMathOperator{\notnsubgp}{\ntriangleleft}
\newcommand{\presentation}[2]{\left<#1\;\middle|\;#2\right>}
\DeclareMathOperator{\Char}{char}
\DeclareMathOperator{\fchar}{char}
\DeclareMathOperator{\triv}{triv}
\DeclareMathOperator{\reg}{reg}
\DeclareMathOperator{\std}{std}
\DeclareMathOperator{\Func}{Func}
\DeclareMathOperator{\End}{End}

% Convex Optimization
\let\sectionSymbol\S
\let\S\relax
\newcommand{\S}{\mathbb{S}}
\DeclareMathOperator{\dist}{dist}
\DeclareMathOperator{\dom}{dom}
\DeclareMathOperator{\diag}{diag}
\DeclareMathOperator{\ones}{\mathbbm{1}}
\newcommand{\minimizeOver}[3]{\begin{array}{rl}\underset{#1}{\ptxt{minimize}} & #2\\ \ptxt{subject to} & #3\end{array}}
\newcommand{\maximizeOver}[3]{\begin{array}{rl}\underset{#1}{\ptxt{maximize}} & #2\\ \ptxt{subject to} & #3\end{array}}
\newcommand{\minimizationProblem}[2]{\minimizeOver{}{#1}{#2}}
\newcommand{\maximizationProblem}[2]{\maximizeOver{}{#1}{#2}}
\newcommand{\minimizeOverUnconstrained}[2]{\begin{array}{rl}\underset{#1}{\ptxt{minimize}} & #2\end{array}}
\newcommand{\maximizeOverUnconstrained}[2]{\begin{array}{rl}\underset{#1}{\ptxt{maximize}} & #2\end{array}}
\newcommand{\minimizationUnconstrained}[1]{\minimizeOverUnconstrained{}{#1}}
\newcommand{\maximizationUnconstrained}[1]{\maximizeOverUnconstrained{}{#1}}
\DeclareMathOperator{\argmin}{argmin}
\DeclareMathOperator{\argmax}{argmax}

% Proofs
\newcommand{\st}{s.t.}
\newcommand{\unique}{!}
\newcommand{\iffdef}{\overset{\ptxt{def}}{\Leftrightarrow}}
\newcommand{\eqVertical}{\rotatebox[origin=c]{90}{=}}
\newcommand{\mapsfrom}{\mathrel{\reflectbox{\ensuremath{\mapsto}}}}
\newcommand{\mapsdown}{\rotatebox[origin=c]{-90}{$\mapsto$}\mkern2mu}
\newcommand{\mapsup}{\rotatebox[origin=c]{90}{$\mapsto$}\mkern2mu}
\newcommand{\from}{\!\mathrel{\reflectbox{\ensuremath{\to}}}}
\newcommand{\labeledeq}[1]{\overset{\mathclap{\ptxt{#1}}}{=}}
\newcommand{\eqdef}{\labeledeq{def}}

% Brackets
\newcommand{\paren}[1]{\left(#1\right)}
\renewcommand{\brack}[1]{\left[#1\right]}
\renewcommand{\brace}[1]{\left\{#1\right\}}
\newcommand{\ang}[1]{\left<#1\right>}

% Algorithms
\algrenewcommand{\algorithmiccomment}[1]{\hskip 1em \texttt{// #1}}
\algrenewcommand\algorithmicrequire{\textbf{Input:}}
\algrenewcommand\algorithmicensure{\textbf{Output:}}
\newcommand{\algP}{\ptxt{\textbf{P}}}
\newcommand{\algNP}{\ptxt{\textbf{NP}}}
\newcommand{\algNPC}{\ptxt{\textbf{NP-Complete}}}
\newcommand{\algNPH}{\ptxt{\textbf{NP-Hard}}}
\newcommand{\algEXP}{\ptxt{\textbf{EXP}}}
\DeclareMathOperator{\fl}{fl}

%%%%%%%%%%%%%%%%%%%%%%%%%%%%%
% Other commands
%%%%%%%%%%%%%%%%%%%%%%%%%%%%%
\newcommand{\flag}[1]{\textbf{\textcolor{red}{#1}}}
\let\uSym\u
\let\u\relax
\newcommand{\u}[1]{\underline{#1}}
\let\bSym\b
\let\b\relax
\newcommand{\b}[1]{\textbf{#1}}
\let\iSym\i
\let\i\relax
\newcommand{\i}[1]{\textit{#1}}
\let\scSym\sc
\let\sc\relax
\newcommand{\sc}[1]{\textsc{#1}}

%%%%%%%%%%%%%%%%%%%%%%%%%%%%%%%%%%%%%%%
% Make l's curvy in math environments %
%%%%%%%%%%%%%%%%%%%%%%%%%%%%%%%%%%%%%%%
\mathcode`l="8000
\begingroup
\makeatletter
\lccode`\~=`\l
\DeclareMathSymbol{\lsb@l}{\mathalpha}{letters}{`l}
\lowercase{\gdef~{\ifnum\the\mathgroup=\m@ne \ell \else \lsb@l \fi}}%
\endgroup

%%%%%%%%%%%%%%%%%%%%%%%%%
% Fix \vdots and \ddots %
%%%%%%%%%%%%%%%%%%%%%%%%%
\usepackage{letltxmacro}
\LetLtxMacro\orgvdots\vdots
\LetLtxMacro\orgddots\ddots

\makeatletter
\DeclareRobustCommand\vdots{%
	\mathpalette\@vdots{}%
}
\newcommand*{\@vdots}[2]{%
	% #1: math style
	% #2: unused
	\sbox0{$#1\cdotp\cdotp\cdotp\m@th$}%
	\sbox2{$#1.\m@th$}%
	\vbox{%
		\dimen@=\wd0 %
		\advance\dimen@ -3\ht2 %
		\kern.5\dimen@
		% remove side bearings
		\dimen@=\wd2 %
		\advance\dimen@ -\ht2 %
		\dimen2=\wd0 %
		\advance\dimen2 -\dimen@
		\vbox to \dimen2{%
			\offinterlineskip
			\copy2 \vfill\copy2 \vfill\copy2 %
		}%
	}%
}
\DeclareRobustCommand\ddots{%
	\mathinner{%
		\mathpalette\@ddots{}%
		\mkern\thinmuskip
	}%
}
\newcommand*{\@ddots}[2]{%
	% #1: math style
	% #2: unused
	\sbox0{$#1\cdotp\cdotp\cdotp\m@th$}%
	\sbox2{$#1.\m@th$}%
	\vbox{%
		\dimen@=\wd0 %
		\advance\dimen@ -3\ht2 %
		\kern.5\dimen@
		% remove side bearings
		\dimen@=\wd2 %
		\advance\dimen@ -\ht2 %
		\dimen2=\wd0 %
		\advance\dimen2 -\dimen@
		\vbox to \dimen2{%
			\offinterlineskip
			\hbox{$#1\mathpunct{.}\m@th$}%
			\vfill
			\hbox{$#1\mathpunct{\kern\wd2}\mathpunct{.}\m@th$}%
			\vfill
			\hbox{$#1\mathpunct{\kern\wd2}\mathpunct{\kern\wd2}\mathpunct{.}\m@th$}%
		}%
	}%
}
\makeatother

\newcommand{\B}{
	\begin{tikzpicture}
	\filldraw [fill=red, draw=black] (0, 0) rectangle (0.37, 0.45);
	\draw [line width=0.5mm, white ] (0.1,0.08) -- (0.1,0.38);
	\draw[line width=0.5mm, white ] (0.1, 0.35) .. controls (0.2, 0.35) and (0.4, 0.2625) .. (0.1, 0.225);
	\draw[line width=0.5mm, white ] (0.1, 0.225) .. controls (0.2, 0.225) and (0.4, 0.1625) .. (0.1, 0.1);
	\end{tikzpicture}
}

% Allow custom symbols as arrows in tikz-cd
\tikzset{
  symbol/.style={
    draw=none,
    every to/.append style={
      edge node={node [sloped, allow upside down, auto=false]{$#1$}}}
  }
}
% Example: \arrow[r,symbol=\cong]
% https://tex.stackexchange.com/questions/394154/how-to-include-inclusion-subgroup-relationship-in-tikz-cd-diagram

\author{Professor Alejandro Uribe-Ahumada\\ \small\i{Transcribed by Thomas Cohn}}
\title{Math 591 Lecture 35}
\date{11/30/20} % Can also use \today

\begin{document}
\maketitle
\setlength\RaggedRightParindent{\parindent}
\RaggedRight

\par\noindent
Final remarks on orientation:\n

\par\noindent
Recall: If $M$ is oriented, $\exists\set{(U_{\alpha},\varphi_{\alpha})}$, a positive atlas of $M$. This means all transition functions $\varphi_{\beta}\of\varphi_{\alpha}\inv$ have the determinant of their Jacobian positive at every point, and the coordinate frames are positive.\n

\par\noindent
Conversely, if $M$ is an atlas satisfying he above property, then one can define an orientation of $M$ by requiring the coordinate frames are positive.\n

\par\noindent
In general, to show a manifold is orientable, exhibit such an atlas.\n

\ex{
	Check: The atlas of $\CP^{n}$ used in homework has this condition, so it is orientable.\n
}

\par\noindent
Observe: If $S\subset{}M$ is a $\codim$-$1$ submanifold, and $M$ is oriented, and there exists a continuous field of normal vectors on $S$,
\[
	S\ni{}p\mapsto\vec{n}_{p}\in{}T_{p}M\quad\ptxt{\st{}}\quad{}T_{p}M=T_{p}S\oplus\R\vec{n}_{p}
\]
then $S$ is orientable, and the convention for its orientation is: a basis $\set{b_{1},\ldots,b{n-1}}$ of $T_{p}S$ is positive iff $\set{\vec{n}_{p},b_{1},\ldots,b_{n-1}}$ is positive w.r.t. $M$.\n
\i{TL;DR, put the normal vector first.}\n

\ex{
	We can embed the Klein bottle in a $\dim$-$3$ manifold $M$ \st{} there exists a continuous $\vec{n}$, but $M$ is non-orientable.\n
}

\subsection*{Partitions of Unity}

\par\noindent
This is a very technical, but very useful tool. We begin with point-set topology.\n

\defn{
	An indexed covering (not necessarily open) $\set{S_{\alpha}}_{\alpha\in{}A}$ of (a manifold) $M$ (doesn't have to be a manifold) (with $S_{\alpha}\subset{}M$) is said to be \u{locally finite} iff every $p\in{}M$ has a neighborhood $U$ \st{} $\set{\alpha\in{}A\mid{}S_{\alpha}\cap{}U\ne\emptyset}$ is finite. That is, every $p$ is in only finitely-many $S_{\alpha}$.\n
}

\thm{
	(Thm 1.15 in Lee) Any topological manifold is \u{paracompact}: every open cover $\set{U_{\alpha}}_{\alpha\in{}A}$ has a countable, locally-finite refinement $\set{V_{i}}_{i\in\N}$. That is, $\forall{}i\in\N$, $V_{i}$ is open, and $\exists\alpha\in{}A$ \st{} $V_{i}\subset{}U_{\alpha}$ and $M$ is covered by $\set{V_{i}}_{i\in\N}$.\nn
	Proof: This proof is long and complex, but it only uses point-set topology. This is the first time we're using the fact that $M$ is second-countable!\n
}

\par\noindent
Observe: If $\mathscr{B}$ is any basis of $M$, the $V_{i}$ can be chosen to be in $\mathscr{B}$.\n

\defn{
	Let $M$ be a smooth manifold. A \u{partition of unity} on $M$ is an indexed family $\set{\chi_{\alpha}}_{\alpha\in{}A}$ of $C^{\infty}$ functions on $M$ \st{}
	\begin{enumerate}[leftmargin=4\parindent, itemsep=0pt, label=(\arabic*)]
		\item $\set{\supp(\chi_{\alpha})}_{\alpha\in{}A}$ is a locally finite cover of $M$.
		\item $\forall{}p\in{}M$, $\sum_{\alpha\in{}A}\chi_{\alpha}(p)=1\in\R$. (Note that this is a finite sum by (1).)
	\end{enumerate}
}

\thm{
	Let $\set{U_{\alpha}}_{\alpha\in{}A}$ be an open cover of $M$. Then $\exists\set{\chi_{\alpha}}_{\alpha\in{}A}$, a partition of unity, that is \u{subordinate} to $\set{U_{\alpha}}_{\alpha\in{}A}$, i.e., $\forall\alpha\in{}A$, $\supp\chi_{\alpha}\subseteq{}U_{\alpha}$.\nn
	Proof: We'll use paracompactness. (It may be easier to start by just thinking of a compact manifold.)\n
	Let $\mathscr{B}$ be the set of normal coordinate balls in $M$; we define $B\subset{}M$ to be a \u{normal coordinate ball} in $M$ iff there's a chart $(U,\phi)$ \st{} $\closure{B}\subset{}U$ and $\phi(B)=B_{r}(0)\subset\R^{n}$, the ball of radius $r$ centered at $0$ in $\R^{n}$, and also $\exists{}r'>r$ \st{} $\closure{B_{r}(0)}\subset{}B_{r'}(0)\subset\phi(U)$.\nn
	We claim that $\mathscr{B}$ is a basis of the topology of $M$. Let $\set{U_{\alpha}}_{\alpha\in{}A}$ be any open cover. Use the theorem on paracompactness: $\exists\set{B_{i}}_{i\in\N}$, a locally-finite refinement, and $\forall{}i\in\N$, $B_{i}$ is a normal coordinate ball. $\forall{}i\in\N$, let
	\[
		\phi_{i}(B_{i})=B_{r_{i}}(0)\subset\closure{B_{r_{i}}(0)}\subset{}B_{r_{i}'}(0)
	\]
	and $H_{i}$ be a function on $\Im(\phi_{i})$ such that $H_{i}:\Im(\phi_{i})\to\R$ is smooth, with
	\begin{itemize}[leftmargin=4\parindent, label=-]
		\item $H_{i}>0$ on $B_{r_{i}}(0)$
		\item $H_{i}=0$ on $B_{r_{i}}(0)\comp$
	\end{itemize}\up\n
	Thus, $\supp{}H_{i}=\closure{B_{r_{i}}(0)}$.\nn
	Define $\psi_{i}\in{}C^{\infty}(M)$ by $\psi_{i}=H_{i}\of\phi_{i}$ on $\dom\phi_{i}$, and $0$ everywhere else. Then $\supp\psi_{i}=\closure{B_{i}}\subset{}M$. We claim that $\set{\closure{B_{i}}}_{i\in\N}$. Observe that $\forall{}p$, $\sum_{i\in\N}\psi_{i}(p)>0$, because $\set{B_{i}}_{i\in\N}$ forms a cover of $M$, and $\restr{\psi_{i}}{B_{i}}>0$.\nn
	Now, define
	\[
		f_{i}=\frac{1}{\displaystyle\sum_{j\in\N}\psi_{j}}\psi_{i}
	\]
	so that $\set{\supp{}f_{i}}_{i\in\N}=\set{\closure{B_{i}}}$ is locally finite, and $\sum_{i\in\N}f_{i}=1$, $\forall{}p\in{}M$. Then, we just have to fix it so that he indexing sets are the same as $\set{U_{\alpha}}_{\alpha\in{}A}$ by $\forall{}i\in\N$, pick $\alpha(i)\in{}A$ such that $B_{i}\subset{}U_{\alpha(i)}$ and $\forall\alpha\in{}A$, let
	\[
		\chi_{\alpha}=\sum_{\substack{i\ptxt{ \st{}}\\ \alpha(i)=\alpha}}f_{i}
	\]
	(Note that $\chi_{\alpha}=0$ if the sum is empty.)\nn
	We claim that $\set{\supp\chi_{\alpha(i)}}_{i\in\N}$ is still locally finite. This follows from $\set{\supp{}f_{i}}_{i\in\N}$ being locally finite.\proven
}

\par\noindent
There are many applications of partitions of unity!

\ex{
	Existence of $C^{\infty}$ sections of any vector bundle.\nn
	Say $E\to{}M$ is a vector bundle of rank $r$. Then there exist $\set{(U_{\alpha},f_{\alpha})}$ local trivializations:
	\[
		\begin{tikzcd}[ampersand replacement=\&]
			\pi\inv(U_{\alpha}) \arrow[rr, "f_{\alpha}"] \arrow[dr] \&[-10pt] \&[-10pt] U_{\alpha}\times\R^{n} \arrow[dl]\\
			\& U_{\alpha} \arrow[ul, bend left=45, "s_{\alpha}"] \arrow[ur, bend right=45, "\ptxt{pick any section $s$}", swap]
		\end{tikzcd}
	\]
	Let $\set{\chi_{\alpha}}$ be a partition of unity on $M$ subordinate to $\set{U_{\alpha}}$. Then let $s=\sum_{\alpha\in{}A}\chi_{\alpha}\cdot{}s_{\alpha}$ (we interpret $\chi_{\alpha}\cdot{}s_{\alpha}$ as a $C^{\infty}$ section on $M$).\n
}

\par\noindent
The main application of partitions of unity is integrating forms.\n

\defn{
	Let $M$ be an oriented $n$-dimensional manifold. Let $\mu\in\Omega^{n}_{0}(M)$ be a top degree form with compact support. Let $\set{\phi_{\alpha}}$ be a positive atlas, and $\set{\chi_{\alpha}}$ a subordinate partition of unity (i.e. $\supp\chi_{\alpha}\subseteq\supp\phi_{\alpha}$, $\forall\alpha$). Then we define
	\[
		\int_{M}\mu=\sum_{\alpha}\int(\phi_{\alpha}\inv)^{*}(\chi_{\alpha}\mu)
	\]
}

\par\noindent
We have to check that the right hand side is independent of choice of coordinates. We'll do this next time...

\end{document}